\section*{第 \ref{ch:g8:reproduction-and-environment} 章——繁殖与环境}

\begin{solution}{exo:g8:reproduction-and-environment:1}
规律是：幼体须在食物与\omterm{def:g6:exploring-our-environment:conditions}{条件}都适宜时出生，
对山雀雏鸟而言，正是栎树上春季\omterm{ex:g2:animals-grow-up:butterfly}{毛毛虫}的高峰。
促使其生殖机制启动的信号是
日照逐渐变长。
\end{solution}

\begin{solution}{exo:g8:reproduction-and-environment:2}
春季产仔，那时青草\omterm{def:g2:from-seed-to-plant:germination}{萌发}、气候温和，
正好迎接小牛。妊娠的月数须从出生日往前倒推
——秋季\omterm{prop:g4:animal-reproduction:where}{交配}，正是这一倒算的结果。
\end{solution}

\begin{solution}{exo:g8:reproduction-and-environment:3}
因为每一枚卵、每一粒\omterm{def:g2:from-seed-to-plant:seed}{种子}和每一滴\omterm{def:g2:animals-grow-up:born}{乳汁}都是\omterm{def:g6:plants-are-producers:matter}{有机物}，
亲代必须先获取并重新组建这些\omterm{def:g6:plants-are-producers:matter}{有机物}，
（见\cref{def:g6:matter-from-food:production}）。代价有：
母鸡每日产的一枚卵、山毛榉大年结的\omterm{def:g2:from-seed-to-plant:seed}{种子}，
母鹿的\omterm{def:g2:animals-grow-up:born}{乳汁}（或公鹿每年再生的鹿角）。
\end{solution}

\begin{solution}{exo:g8:reproduction-and-environment:4}
同一地点内同一\omterm{def:g5:biodiversity:species}{物种}的全部个体。
繁殖（与迁入）使其增加，死亡与迁出使其减少；
\omterm{def:g6:exploring-our-environment:components}{环境}中的食物、空间、捕食者，
以及各种\omterm{def:g6:exploring-our-environment:conditions}{条件}，共同设定其数量上限。
\end{solution}

\begin{solution}{exo:g8:reproduction-and-environment:5}
捕食（苍鹭、蜻蜓幼虫）；食物与空间
（食物有限，栖身之所也有限）；\omterm{def:g6:exploring-our-environment:conditions}{条件}
（霜冻、干旱、池塘萎缩）。
\end{solution}

\begin{solution}{exo:g8:reproduction-and-environment:6}
\omterm{def:g2:where-animals-live:habitat}{栖息地}的质量写在卵的数目上：
毛虫丰富的\omterm{def:g2:where-animals-live:habitat}{栖息地}，能养出可哺育十只雏鸟的母鸡；
贫瘠的则只能哺育六到七只。
\end{solution}

\begin{solution}{exo:g8:reproduction-and-environment:7}
母羊的繁殖机制由日照渐短启动；
人工光照按饲养者的安排重放这一信号，
羊群因此提前产羔。这证明繁殖时机是
直接从信号本身读取的——改变光照，
“日历”便随之改变。
\end{solution}

\begin{solution}{exo:g8:reproduction-and-environment:8}
时机：信号与产卵都正常——没有错配。
投入：亲代供给正常，产下数以千计的
卵。检查：三项全部收紧——干旱减少了空间与\omterm{def:g7:what-is-respiration:respiration}{氧气}，
拥挤造成食物匮乏，苍鹭集中捕猎，
\omterm{def:g2:where-animals-live:habitat}{栖息地}：繁殖池本身几近消失——
这是决定性的一条。
\end{solution}

\begin{solution}{exo:g8:reproduction-and-environment:9}
雏鸟孵出时，盛宴早已结束：\omterm{ex:g2:animals-grow-up:butterfly}{毛毛虫}
高峰已过，而日照信号仍在提示“再等等”。
出飞数量仍在下降——天气、投入与捕食状况都还有利，
但时机这一项未过审核：信号与资源
的高峰已不再同步。
\end{solution}

\begin{solution}{exo:g8:reproduction-and-environment:10}
数以百万计的卵正是这一策略：每枚卵都是押上的一注，
要与捕食、饥饿与\omterm{def:g6:exploring-our-environment:conditions}{条件}抗衡，检查
几乎将其全部收走——这正是计划本身，
而非计划的失败。平均而言，存活者取代亲代；
投入\emph{正是}不计后果的策略
换取少数赢家的方式（见\cref{exo:g4:animal-reproduction:7}，
此点在那里最早提出）。
\end{solution}

\begin{solution}{exo:g8:reproduction-and-environment:11}
排干沼泽：移除繁殖\omterm{def:g2:where-animals-live:habitat}{栖息地}——
即\omterm{def:g2:where-animals-live:habitat}{栖息地}一项。巢箱：增加空间，缓解了那一项
检查。五月喷药清空了食物高峰——正是食物那一项，
且恰在最糟糕的时刻。\omterm{def:g7:organs-at-work:recovery}{恢复}湿草地：
同时\omterm{def:g7:organs-at-work:recovery}{恢复}\omterm{def:g2:where-animals-live:habitat}{栖息地}与食物，
（见\cref{ex:g5:biodiversity:storks}）。
\end{solution}

\begin{solution}{exo:g8:reproduction-and-environment:12}
青蛙处在策略线“多产而不管”的一端：
成千上万枚廉价的卵，无人守护。金字塔的算术
保证大部分\omterm{ex:g2:animals-grow-up:frog}{蝌蚪}的\omterm{def:g6:plants-are-producers:matter}{有机物}必须用于养活
池塘中的\omterm{def:g5:food-webs:roles}{消费者}，而非变成青蛙——每一级
只能传递其中一小部分。而正是这些检查
使这一算术得以兑现：捕食、食物、空间与
\omterm{def:g6:exploring-our-environment:conditions}{条件}每季收走过剩的部分，平均而言，
只留下一对来替代一对。池塘之所以未被青蛙挤满，
是因为策略已为损耗作了安排，金字塔
为之定价，检查则将其收走。
\end{solution}

\begin{solution}{pb:g8:reproduction-and-environment:1}
\textbf{1.} 春季渐长的日照启动产卵；
时机的安排使饥饿的雏鸟赶上栎树上的
\omterm{ex:g2:animals-grow-up:butterfly}{毛毛虫}高峰。

\textbf{2.} \Cref{prop:g8:reproduction-and-environment:cost}：
繁殖的供给来自\omterm{def:g2:where-animals-live:habitat}{栖息地}——
在丰饶季节中积累起来的亲代，之后能产下更大规模的卵。

\textbf{3.} 卵与雏鸟遭到捕食；寒潮中食物短缺；
\omterm{def:g6:exploring-our-environment:conditions}{条件}方面，巢中一周阴冷潮湿的天气；
（空间则由巢箱本身提供）。

\textbf{4.} 因为卵数本身什么也说明不了：核查需要
提示与高峰的时间安排（天气记录）、预算（季节所提供的）、
检查以及\omterm{def:g2:where-animals-live:habitat}{栖息地}（园地工作）。
账本只有与
其背景栏并列才可读。

\textbf{5.} 雏鸟在早期高峰已过之后才孵出——
这就是时间线：提示未变，
资源转移，巧合被打破。

\textbf{6.} 预算与检查经由\omterm{def:g2:where-animals-live:habitat}{栖息地}体现：枯枝
承载\omterm{prop:g3:sorting-living-things:groups}{昆虫}与天然洞穴，每月刈割的草带
既不产\omterm{ex:g2:animals-grow-up:butterfly}{毛毛虫}也不产蜘蛛——食物检查收紧，
\omterm{def:g2:where-animals-live:habitat}{栖息地}贫瘠化，
尽管天气宜人。

\textbf{7.} 捕食。卵的供给在鹰的
捕食到来之前已经完成——由雌鸟的预算决定；
鹰稍后在雏鸟身上取走。
账本的不同栏目，收取者也不同。

\textbf{8.} 巢位核查——更多繁殖位点，
因而繁殖对总数增加。未受影响的是
每块领地的食物上限，这也是为什么
每只巢箱自身的成功率保持不变，而非上升。

\textbf{9.} “十年的窝卵数追随园地：富足的\omterm{ex:g2:animals-grow-up:butterfly}{毛毛虫}
年份与杂乱的栎树，使巢箱里装满八至十枚卵；
贫瘠且整饬的年份只有六枚——
窝卵数即\omterm{def:g2:where-animals-live:habitat}{栖息地}质量，以卵的数目写出。”

\textbf{10.} 巢箱只替代巢位空间。栎树同时
是提示的目标与预算——时间安排所指向的
\omterm{ex:g2:animals-grow-up:butterfly}{毛毛虫}高峰，以及供养窝卵与雏鸟的食物——
另外还是减轻检查的因素（其枝条天然的储藏）。
停车场上的五十个巢箱，
不过是空了的食品柜里的几层架子。

\textbf{11.} 春季期间让草地带不予刈割——
补充食物预算；保留（或补植）枯枝与老栎树——
保护\omterm{def:g2:where-animals-live:habitat}{栖息地}，以及
整个日历所指向的高峰。

\textbf{12.} 一个\omterm{def:g2:where-animals-live:habitat}{栖息地}要有：适时的提示、
支付提示所需的食物，以及
足够温和的检查，使平均每对亲鸟都能留下一对接替者。
\end{solution}
